\setcounter{section}{1}

\Section{Теорема Кантора о вложенных отрезках}

\begin{definition}
    \textit{Системой вложенных отрезков} называют последовательность отрезков
    \([a_{n}, b_{n}]\), в которой каждый следующий отрезок содержится
    в предыдущем:
    \[[a_{n+1}, b_{n+1}] \subset [a_{n}, b_{n}],\qquad n\in\N.\]
\end{definition}

\begin{definition}
    \textit{Системой стягивающихся отрезков} называют систему вложенных отрезков,
    для которой длина отрезка стремится к нулю:
    \[\lim_{n\to\infty}(b_{n}-a_{n})=0.\]
\end{definition}

\begin{theorem}[Теорема об ограниченной монотонной последовательности]
    Если последовательность ограничена и монотонна, то она сходится.
\end{theorem}

\begin{theorem}[Теорема о трех последовательностях]
    Если последовательности \(a_{n}\) и \(b_{n}\) сходятся к одному и тому же числу,
    а для последовательности \(c_{n}\) справедливо
    \[\forall\,n\in\N \qquad a_{n} \leq c_{n} \leq b_{n}\]
    То последовательность \(c_{n}\) также сходится, причем
    \[\lim_{ n \to \infty } a_{n} = \lim_{ n \to \infty } c_{n} = \lim_{ n \to \infty } b_{n}.\]
\end{theorem}

    \textit{Примечание: в рамках сдачи этой темы теоремы не требуют доказательства.
    Доказательство теоремы 2.1 смотри в следующей теме.}

    \vspace{2mm}

\begin{theorem}[Теорема Кантора]
    Для любой системы вложенных отрезков существует точка, принадлежащая всем отрезкам системы. 
    Если система стягивается, то эта точка - единственна.

    \vspace{2mm}

    \proofname~Рассмотрим произвольную систему вложенных отрезков.
    Из вложенности следует, что последовательность \(a_{n}\) не убывает,
    а последовательность \(b_{n}\) не возрастает. Кроме того,
    \(a_{n}\leq b_{n}\) для любого \(n\). Поэтому последовательность
    \(a_{n}\) ограничена сверху, а последовательность \(b_{n}\) ограничена
    снизу. Следовательно, по теореме об ограниченной монотонной 
    последовательности существуют пределы
    \[a=\lim_{n\to\infty}a_{n},\qquad
    b=\lim_{n\to\infty}b_{n}.\]
    Так как \(a_{n}\leq b_{n}\) при любом \(n\), то по свойству пределов \(a\leq b\).
    Рассмотрим все возможные случаи:

    \begin{itemize}
        \item Если \(a=b\), обозначим это число через \(c\). Поскольку
        \(a_{n}\leq a\) и \(b\leq b_{n}\), получаем
        \[a_{n}\leq c\leq b_{n}\quad\text{для любого }n.\]
        Значит, точка \(c\) принадлежит каждому отрезку системы.

        \item Если \(a<b\), то для любого \(n\) имеем
        \[a_{n}\leq a<b\leq b_{n}.\]
        Следовательно, отрезок \([a,b]\) содержится в каждом отрезке
        системы, и любая точка из этого отрезка удовлетворяет теореме.
    \end{itemize}

    Теперь рассмотрим систему стягивающихся отрезков. Из равенств
    \(\underset{n\to\infty}{\lim}a_{n} = a\) и 
    \(\underset{n\to\infty}{\lim}b_{n} = b\) следует 
    \[b-a=\lim_{ n \to \infty }b_{n} - \lim_{ n \to \infty } a_{n} =
    \lim_{n\to\infty}(b_{n}-a_{n})=0.\]
    Поэтому \(a=b\). Обозначим их общее значение через \(c\). Как было
    показано выше, \(a_{n}\leq c\leq b_{n}\) для любого \(n\), то есть
    точка \(c\) принадлежит всем отрезкам.

    Докажем единственность. Пусть точка \(c'\) также принадлежит каждому
    отрезку системы. Тогда \(a_{n}\leq c'\leq b_{n}\) для любого \(n\).
    Переходя к пределу, по теореме о трех последовательностях
    получаем \(c\leq c'\leq c\), откуда \(c'=c\).
    Следовательно, общая точка единственна.~\qedsymbol
\end{theorem}
